% Source: https://artofproblemsolving.com/community/c3h3176269_2023_mbmt_geometry_round__montgomery_blair_math_tournament
% Stable identity: aops-topic-3176269-q11
Take $O=(0,0)$, $A=(10,0)$, and choose the orientation with $B=(-\dfrac{14}{5},\dfrac{48}{5})$; these coordinates follow from $OB=10$ and $AB=16$. Since $PO=PA$, every possible $P$ lies on $x=5$. The distance from $B$ to this line is $5+\dfrac{14}{5}=\dfrac{39}{5}$, so $PB$ is at least this value. The perpendicular foot $P=(5,\dfrac{48}{5})$ is attainable: the line $P\cdot C=50$ intersects the radius-$10$ circle at $A$ and a distinct second point $C$, since $|P|>5$. The two intersections are not antipodal because the line does not pass through $O$. Thus $AOC$ is nondegenerate and the minimum is $\boxed{\dfrac{39}{5}}$.
\begin{asy}
usepackage("amsmath");
size(7cm);
pair O=(0,0),A=(10,0),Bp=(-14/5,48/5),P=(5,48/5);draw(circle(O,10));draw((5,-11)--(5,13),dashed);draw(O--A^^A--P--O^^Bp--P,blue);label("$O$",O,SW);label("$A$",A,E);label("$B$",Bp,NW);label("$P$",P,NE);
\end{asy}
